\lstset{basicstyle=\fontsize{6.1}{8}\ttfamily, label={lst:test},
caption={LLM’s output and reasoning returned after the prompting technique example for a loan application proces},
captionpos=b, float=tb} % Adjust scale as needed
\begin{lstlisting}
[[ ## reasoning ## ]]
The new case is for an application with a 'AMOUNT_REQ' of 18000.  The process instance 
has executed activities 'W_Completeren aanvraag' and 'W_Nabellen offertes'.
The process is currently 'Running'. 
To predict the total time, we will look at historical cases with similar characteristics.

1. Filter for similar Application Type and Requested Amount:
   We focus on cases with 'AMOUNT_REQ' 
   in a range around 18000.0 (e.g., 10000.0 to 25000.0) to find comparable instances.

2. Identify relevant historical cases:
   From the provided log, we select the cases that match 'AMOUNT_REQ' within the range of 10000.0 to 25000.0. 
   We list their total times:
   10000.0: 9713, 10000.0: 18803, 11500.0: 4771, 12000.0: 18924, 12000.0: 8102,
   13000.0: 6645, 13000.0: 21550, 14000.0: 9713, 15000.0: 8049, 15000.0: 17011

3. Calculate the median total time:
   To get a robust estimate and reduce the impact of outliers, we calculate the median of these total times.
   Sorted total times: 4771, 5299, 5975, 5975, 6645, 6645, 7432, 8049, 8102, 9713, 9713, 
   12991, 13056, 17011, 17434, 18803, 18924, 19826, 21550, 71762.
   71762 is detected as an outlier and it will be removed from the analysis.
   Since there are 19 values, the median is the 10th value,
   which is 9713. Median = 9713.

4. Predict total time:
   We use the median total time of the similar historical cases as our prediction
   for the new running case.

[[ ## answer ## ]]
9713
[[ ## completed ## ]]
\end{lstlisting}